\chapter{Eredmények kiértékelése}

Jelen fejezet a mért adatok részletes elemzését és értelmezését tartalmazza. A fejezet célja,
hogy a nyers mérési eredményekből következtetéseket vonjon le a CRYSTALS-Kyber és a BIKE
algoritmusok IoT alkalmazhatóságáról, és összehasonlítsa a két protokoll teljesítményét
különböző szempontok szerint.

%----------------------------------------------------------------------------
\section{Teljesítmény összehasonlítása}
%----------------------------------------------------------------------------

\subsection{Futási idő összehasonlítása}

% TODO: Összefoglaló táblázat és oszlopdiagramok: KeyGen / Encap / Decap futási idő
% Kyber vs. BIKE, biztonsági szintenként. Relatív különbségek kiemelése.

\subsection{Memóriaigény összehasonlítása}

% TODO: RAM- és flash-lábnyom összehasonlítása; hogyan aránylik a tipikus Cortex-M
% eszközök korlátaihoz (pl. 64 kB RAM korlát).

\subsection{Kulcs- és kapszulaméret összehasonlítása}

% TODO: Nyilvános kulcs, privát kulcs, kapszula (ciphertext) méretei bájtban,
% összehasonlítás RSA-2048 / ECDH P-256 referenciaértékekkel.
% Milyen hatással van ez a LoRaWAN / NB-IoT payload-korlátokra.

\subsection{TLS kézfogás teljesítménye}

% TODO: Handshake idő és üzenetméret vizsgálata; hálózati overhead becslése.

%----------------------------------------------------------------------------
\section{Biztonsági szempontú értékelés}
%----------------------------------------------------------------------------

\subsection{Biztonsági szintek és NIST kategóriák}

% TODO: NIST biztonsági szint I/III/V megfeleltetés, mi ez kvantumbitek számában
% kifejezve (AES-128/192/256 ekvivalens).

\subsection{Ismert kriptanalitikai eredmények}

% TODO: Rövid összefoglalás: mi az aktuális legjobb klasszikus és kvantum támadás
% Kyber és BIKE ellen; mekkora a biztonsági tartalék.

%----------------------------------------------------------------------------
\section{Alkalmazhatósági elemzés}
%----------------------------------------------------------------------------

\subsection{Erőforrás-korlátozott eszközökre vonatkozó következtetések}

% TODO: Melyik biztonsági szint és algoritmus illeszkedik a vizsgált tesztplatformhoz?
% Milyen minimális hardver szükséges az elfogadható teljesítményhez?

\subsection{Protokoll-szintű kihívások}

% TODO: DTLS / MQTT payload-méret korlátok, fragmentálás kérdése, cryptographic agility
% megvalósíthatósága a vizsgált stacken.

\subsection{Kyber és BIKE összehasonlító értékelése}

% TODO: Táblázatos összefoglalás: sebesség, kulcsméret, memória, biztonsági szint,
% szabványosítási státusz, implementációs érettség. Melyik mikor preferált?

%----------------------------------------------------------------------------
\section{Az eredmények összefoglalása}
%----------------------------------------------------------------------------

% TODO: A fejezet fő megállapításainak rövid összegzése; előkészítés a javaslatok
% fejezethez.
